\subsection{Case Study}
\label{sec:casestudy}

% Case study: pre-registered session-selection protocol and the position of the
% selected sessions within their own condition.
%
% Implementation source: private ML repository commit
% 99aae0af8c3fa1ceb784083446e83c40d0fb917f (2026-08-04).
% Workbooks (implementation-repository relative):
%   data/0803_FULL/SM-S948U_metrics_12MP_normal_0803.xlsx
%   data/0803_FULL/SM-S948U_metrics_24MP_memory_0803.xlsx
%
% Both panels are produced by scripts/export_casestudy.py, which fails unless
% the criteria leave exactly one session per condition.  Per-shot traces are in
% data/casestudy/*.csv; the peer statistics are in
% data/casestudy/*_peer_comparison.csv.

\begin{table}[t]
  \centering
  \setlength{\abovecaptionskip}{2pt}%
  \setlength{\belowcaptionskip}{3pt}%
  \caption{Case-study session selection. The criteria are applied before
  inspecting outcomes and leave exactly one session in each condition; the
  12MP session is the one shown in Fig.~\ref{fig:casestudy_12mp}.}
  \label{tab:casestudy_selection}

  {\scriptsize
  \renewcommand{\arraystretch}{1.06}%
  \setlength{\tabcolsep}{3.5pt}%
  \begin{tabular}{|l||r|r|}
    \hline\hline
    \multicolumn{1}{|c||}{\multirow{2}{*}{\textbf{Selection step}}} &
    \multicolumn{2}{c|}{\textbf{Sessions remaining}} \\
    \cline{2-3}
    & \multicolumn{1}{c|}{\makecell[c]{\textbf{12MP}\\\textbf{normal}}}
    & \multicolumn{1}{c|}{\makecell[c]{\textbf{24MP mode +}\\\textbf{memory}}} \\
    \hline\hline
    Recorded sessions & 44 & 44 \\
    \hline
    C1~~30 captures, no Capture Timeout & 43 & 38 \\
    \hline
    C2~~pacing precedes the first demotion & 28 & 20 \\
    \hline
    C3~~\(M\) and \(S\) demoted at distinct captures & 4 & 4 \\
    \hline
    C4~~pacing relaxes to \(0\) for \(\ge 2\) captures & 1 & 2 \\
    \hline
    C5~~pacing reactivates before the \(S\) demotion & 1 & 2 \\
    \hline
    C6~~no Lv5 MP12 resolution fallback\rlap{$^{\dagger}$} & 1 & 1 \\
    \hline\hline
    \textbf{Selected session} &
    \makecell[r]{Run 30\\(Lv4)} & \makecell[r]{Run 34\\(Lv3)} \\
    \hline\hline
    \multicolumn{3}{c}{} \\[-7pt]
    \hline\hline
    \multicolumn{1}{|c||}{\multirow{2}{*}{\textbf{Outcome vs.\ same-condition sessions}}} &
    \multicolumn{2}{c|}{\textbf{Selected / median [min--max]}} \\
    \cline{2-3}
    & \multicolumn{1}{c|}{\(n_p=10\)} & \multicolumn{1}{c|}{\(n_p=6\)} \\
    \hline\hline
    Cumulative applied delay (s) &
      \makecell[r]{5.21 / 4.41\\{}[2.93--8.38]} &
      \makecell[r]{4.39 / 4.23\\{}[2.69--6.62]} \\
    \hline
    Paced transitions (\%) &
      \makecell[r]{55.2 / 43.1\\{}[20.7--75.9]} &
      \makecell[r]{31.0 / 39.7\\{}[13.8--58.6]} \\
    \hline
    Deadline margin \(P_5\) (ms) &
      \makecell[r]{366 / 429\\{}[209--718]} &
      \makecell[r]{302 / 332\\{}[157--769]} \\
    \hline
    Burst span (s) &
      \makecell[r]{21.7 / 20.8\\{}[16.9--22.5]} &
      \makecell[r]{22.8 / 22.7\\{}[19.7--27.3]} \\
    \hline
    \(M\)~executed (\%) &
      \makecell[r]{26.7 / 38.3\\{}[23.3--100]} &
      \makecell[r]{16.7 / 33.3\\{}[10.0--73.3]} \\
    \hline
    \(S\)~executed (\%) &
      \makecell[r]{73.3 / 100\\{}[73.3--100]} &
      \makecell[r]{53.3 / 100\\{}[53.3--100]} \\
    \hline\hline
  \end{tabular}}

  \par\vspace{2pt}
  \begin{minipage}{0.98\columnwidth}
  {\scriptsize \(^{\dagger}\)Binding only in 24MP mode, where Lv5--6 sessions
  use the production MP12 fallback; the 12MP condition has no such fallback.
  Peer sets hold the condition, the starting overheat level, and the capture
  size bucket of the selected session fixed, keep only complete timeout-free
  sessions, and include the selected session itself. In both conditions the
  selected session pays more cumulative pacing delay, holds a smaller
  fifth-percentile deadline margin, spans a longer burst, and completes less
  optional work than its peer median, and it sits at the bottom of its peer set
  for Filter completion; the only metric on the favourable side of the median
  is the paced-transition share of the 24MP session. The sessions were selected
  for mechanism coverage, not for performance. The 24MP-mode column is retained
  to show that the same criteria, applied to a different capture condition,
  again converge on a single session that reproduces the same coordination
  sequence (pacing at capture~4, \(M\) demoted at~6, pacing relaxed over
  captures~6--8, reactivated at~9, \(S\) demoted at~17, no Capture Timeout).
  \par}
  \end{minipage}
\end{table}

% Notes: docs/exhibits.md#fig_casestudy_12mp
\begin{figure}[t]
  \centering
  \setlength{\abovecaptionskip}{2pt}%
  \setlength{\belowcaptionskip}{2pt}%
  \providecommand{\caseStudyDataDir}{data/case_study}
  \providecommand{\csTimeoutRange}[2]{%
    \fill[cs tmo range] ({axis cs:#1,0}|-{rel axis cs:0,0})
      rectangle ({axis cs:#2,0}|-{rel axis cs:0,1});}
  \tikzset{
    cs tmo range/.style={red!70!black,opacity=0.10},
  }
  \pgfplotsset{
    cs executed/.style={
      only marks,mark=square*,mark size=1.35pt,
      draw=black,fill=black!72,line width=0.4pt,
    },
    cs skipped/.style={
      only marks,mark=square,mark size=1.35pt,
      draw=black!45,line width=0.4pt,
    },
    cs delay/.style={
      ybar,bar width=2.1pt,
      draw=blue!55!black,fill=blue!55!black,fill opacity=0.55,line width=0.35pt,
    },
    cs backlog/.style={
      line width=0.7pt,draw=blue!55!black,
      mark=*,mark size=0.9pt,
      mark options={solid,draw=none,fill=blue!55!black},
    },
    cs slack/.style={
      line width=0.7pt,draw=black!80,
      mark=*,mark size=0.9pt,
      mark options={solid,draw=none,fill=black!80},
    },
    cs floor/.style={red!70!black,line width=0.6pt,densely dashed},
    cs step/.style={
      const plot,no marks,line width=0.65pt,draw=black!80,
    },
  }

  \begin{tikzpicture}
    \begin{groupplot}[
      group style={
        group size=1 by 4,
        vertical sep=3.5pt,
        xticklabels at=edge bottom,
      },
      scale only axis=true,
      width=0.78\columnwidth,
      xmin=0.5,xmax=30.5,
      xtick={5,10,15,20,25,30},
      tick label style={font=\scriptsize},
      label style={font=\scriptsize},
      grid=major,
      major grid style={gray!16,line width=0.2pt},
      axis line style={line width=0.4pt},
      tick style={line width=0.4pt},
      tick align=outside,
      tick pos=left,
      ylabel style={
        at={(axis description cs:-0.15,0.5)},
        anchor=south,
        align=center,
      },
    ]

      \nextgroupplot[
        height=0.115\columnwidth,
        ymin=0.4,ymax=2.6,
        ytick={1,2},
        yticklabels={{\(S\)},{\(M\)}},
        yticklabel style={font=\tiny},
        grid=none,
        clip=false,
      ]
      \csTimeoutRange{8}{12}
      \addplot[cs executed] table[x=shot,y=lane,col sep=comma]
        {\caseStudyDataDir/12mp_normal_bokeh_exec.csv};
      \addplot[cs skipped] table[x=shot,y=lane,col sep=comma]
        {\caseStudyDataDir/12mp_normal_bokeh_skip.csv};
      \addplot[cs executed] table[x=shot,y=lane,col sep=comma]
        {\caseStudyDataDir/12mp_normal_filter_exec.csv};
      \addplot[cs skipped] table[x=shot,y=lane,col sep=comma]
        {\caseStudyDataDir/12mp_normal_filter_skip.csv};
      \node[anchor=south,font=\tiny,align=center,
            text=red!70!black,inner sep=1pt]
        at ([yshift=1pt]axis cs:10,2.6)
        {baseline\\first-timeout window};

      \nextgroupplot[
        height=0.135\columnwidth,
        ylabel={Draft Sequence\\queue depth},
        ymin=0,ymax=6.8,
        ytick={0,3,6},
      ]
      \csTimeoutRange{8}{12}
      \addplot[cs step] table[x=shot,y=queue_depth,col sep=comma]
        {\caseStudyDataDir/12mp_normal_backlog.csv};

      \nextgroupplot[
        height=0.22\columnwidth,
        ylabel={Applied delay\\(ms)},
        ymin=0,ymax=1100,
        ytick={0,500,1000},
      ]
      \csTimeoutRange{8}{12}
      \addplot[cs delay] table[x=shot,y=delay_ms,col sep=comma]
        {\caseStudyDataDir/12mp_normal_delay.csv};

      \nextgroupplot[
        height=0.40\columnwidth,
        xlabel={Capture index within the 30-shot run},
        ylabel={\% of\\Capture Timeout},
        ymin=-11,ymax=80,
        ytick={0,20,40,60},
        yticklabels={0\%,20\%,40\%,60\%},
        legend style={
          at={(0.995,0.995)},anchor=north east,
          font=\tiny,draw=none,
          fill=white,fill opacity=0.75,text opacity=1,
          inner xsep=2pt,inner ysep=1pt,
          legend columns=2,
          /tikz/every even column/.append style={column sep=4pt},
        },
        legend cell align=left,
      ]
      \csTimeoutRange{8}{12}
      \addplot[cs floor,forget plot] coordinates {(0.5,0) (30.5,0)};
      \node[anchor=north east,font=\tiny,text=red!70!black,inner sep=1.5pt]
        at ([xshift=-1pt,yshift=-1pt]axis cs:30.5,0)
        {Capture Timeout (Slack \(=0\%\))};
      \addplot[cs backlog] table[x=shot,y expr={\thisrow{backlog_ms}/70},col sep=comma]
        {\caseStudyDataDir/12mp_normal_backlog.csv};
      \addlegendentry{Backlog}
      \addplot[cs slack] table[x=shot,y expr={\thisrow{margin_ms}/70},col sep=comma]
        {\caseStudyDataDir/12mp_normal_margin.csv};
      \addlegendentry{Slack}
      \node[anchor=east,font=\tiny,inner sep=1.5pt] at (axis cs:7.7,2.5)
        {2.1\%};
      \node[anchor=west,font=\tiny,inner sep=1.5pt] at (axis cs:24.4,5.1)
        {5.1\%};
    \end{groupplot}

    \matrix[
      matrix of nodes,
      nodes={font=\tiny,anchor=west,inner ysep=0.4pt},
      column sep=2pt,row sep=0.6pt,anchor=south east,
    ] at ([xshift=5pt,yshift=-2pt]group c1r1.north east) {
      \node[rectangle,draw=black,fill=black!72,line width=0.4pt,
            minimum size=2.7pt,inner sep=0pt] at (0.12,0) {}; & stage executed \\
      \node[rectangle,draw=black!45,line width=0.4pt,
            minimum size=2.7pt,inner sep=0pt] at (0.12,0) {}; & stage skipped by admission \\
    };
  \end{tikzpicture}
  \caption{Case study: stage execution, applied delay, backlog, and slack.}
  \label{fig:casestudy_12mp}
\end{figure}


Table~\ref{tab:casestudy_selection} summarizes the case-study setting and compares the run with its peers.
We selected this run to examine pacing activation, relaxation, and reactivation alongside admission skips.
Its stage-execution counts and slack \(P_5\) are below the peer medians, while both pacing costs are higher.
Figure~\ref{fig:casestudy_12mp} traces these interactions over 30 captures.

\smallskip\noindent\textbf{Control interaction.}
The trace shows a clear division of labor between the two controls.
Admission persistently reduced draft work, while pacing selectively delayed subsequent captures under queue pressure.

Pacing began at capture 5, before slack became critical.
Backlog continued to rise more slowly, but slack still narrowed to 2.1\% of the timeout budget by capture 8.
Admission then began skipping an optional stage at capture 9.
Backlog declined from capture 10, and slack recovered to 27.3\% by capture 15.
Pacing was inactive over captures 13--20 while the stage remained skipped.
Thus, the controller reduced persistent draft work without continuing to add capture delay.

When backlog grew again, pacing resumed at capture 21 and admission began skipping another optional stage at capture 23.
Backlog then declined from its 63.8\% peak, and slack reached 45.6\% by capture 30.
The repeated pattern shows complementary control at two time scales: admission provides persistent workload reduction, while pacing provides reversible timing control as queue pressure changes.